% =====================================================================
%  TEMPLATE MATEMATICA -- IIS "David Maria Turoldo", Zogno (BG)
%  Classe: exam  |  Compilazione: pdflatex (due passaggi)
%  Soluzioni:        decommentare \printanswers
%  Versione BES/DSA: impostare \versioneBEStrue
%  Struttura:        \iniziofila{A} \begin{esercizi} ... \end{esercizi} \finefila
%                    \iniziofila{B} \begin{esercizi} ... \end{esercizi} \finefila
%  Ogni esercizio:   \begin{esercizio}
	%                      \obiettivo{...}
	%                      \question[punti]  (oppure \question + parts con \part[punti])
	%                      testo ...
	%                    \end{esercizio}
%  Nessuno spazio per le risposte: gli studenti scrivono sul foglio a protocollo.
% =====================================================================
\documentclass[11pt,a4paper,addpoints]{exam}

% ---------- Lingua e codifica ----------
\usepackage[utf8]{inputenc}
\usepackage[T1]{fontenc}
\usepackage[italian]{babel}
\usepackage{lmodern}
\usepackage{microtype}

% ---------- Matematica, grafica, tabelle ----------
\usepackage{amsmath,amssymb}
\usepackage{siunitx}
\sisetup{output-decimal-marker={,}, per-mode=symbol, group-separator={\,}}
\usepackage{tikz}
\usetikzlibrary{arrows.meta,calc,patterns,angles,quotes}
\usepackage{pgfplots}
\pgfplotsset{compat=1.18}
\usepackage{graphicx}
\usepackage{multicol}
\usepackage{array,tabularx,booktabs}
\usepackage{xcolor}

% ---------- Pagina ----------
\usepackage[a4paper,top=2cm,bottom=2.2cm,left=2cm,right=2.2cm,
headheight=30pt,headsep=12pt]{geometry}

% ---------- Opzioni ----------
%\printanswers                 % <-- decommentare per la versione con soluzioni
\newif\ifversioneBES
\versioneBESfalse              % <-- \versioneBEStrue per la versione BES/DSA
\newif\iftabellabonus
\tabellabonusfalse          % <-- true se c'e' un esercizio facoltativo (\bonusquestion): aggiunge la riga Bonus

% ---------- Interlinea ----------
%  \interlinea      : testo normale
%  \interlineaparts : sottopunti a), b), ... (piu' ariosa, separa le equazioni)
\newcommand{\interlinea}{1.3}
\newcommand{\interlineaparts}{1.3}
\ifversioneBES
\usepackage[scaled=1.05]{helvet}
\renewcommand{\familydefault}{\sfdefault}
\renewcommand{\interlinea}{1.5}
\renewcommand{\interlineaparts}{1.7}
\fi
%\linespread{1.5}

% ---------- Dati della prova (da compilare) ----------
\newcommand{\materia}{Matematica}
\newcommand{\classe}{2 BS}
\newcommand{\indirizzo}{Liceo Scientifico}
%   Liceo Scientifico | Liceo Scientifico a indirizzo sportivo | Liceo delle Scienze Umane
\newcommand{\tipoprova}{Verifica scritta}
\newcommand{\argomento}{Equazioni di primo grado e scomposizione}
\newcommand{\durata}{50 minuti}
\newcommand{\annoscolastico}{2026/2027}
\newcommand{\docente}{Prof. Rubens Longhi}
\newcommand{\fila}{A}
\newcommand{\indicazioni}{Svolgi gli esercizi sul foglio a protocollo motivando con chiarezza i passaggi algebrici.
	Non \`e consentito l'uso della calcolatrice\ifversioneBES, salvo quanto previsto dal PDP. Sono consentiti gli strumenti compensativi previsti dal PDP\fi.}

% ---------- Intestazione e pie' di pagina ----------
%  Ogni fila ricomincia da pagina 1: la sua prima pagina ha l'intestazione
%  della scuola e la numerazione "Pagina x di y" e' calcolata sulla singola fila.
\pagestyle{headandfoot}
\firstpageheader
{\scriptsize\textsc{IIS ``David Maria Turoldo''}\\[-1pt]
	\scriptsize Via Ronco 11 -- 24019 Zogno (BG)}
{}
{\scriptsize A.S.\ \annoscolastico\\[-1pt]\scriptsize \docente}
\firstpageheadrule
\runningheader{\small\materia\ -- \classe}{\small\tipoprova\ -- Fila \fila}{}
\runningheadrule
\newcommand{\paginefila}{\pageref{fine@fila@\fila}}
\firstpagefooter{}{\small Pagina \thepage\ di \paginefila}{}
\runningfooter{}{\small Pagina \thepage\ di \paginefila}{}

% ---------- Formato degli esercizi ----------
%  Riga 1:  "1."  in grassetto + obiettivo in grigio ........ (n punti)
%  Riga 2+: testo dell'esercizio
% ---------- Formato degli esercizi ----------
\newcommand{\obiettivocorrente}{}
\newcommand{\obiettivo}[1]{%
	\gdef\obiettivocorrente{#1}}
\newcommand{\stampaobiettivo}{%
	\ifx\obiettivocorrente\empty\else
	{\footnotesize\color{black!55}\textit{Obiettivo: \obiettivocorrente}}%
	\fi}
\qformat{\textbf{\thequestion.}\enspace\stampaobiettivo\hfill\colonnapunti{(\totalpoints\ \points)}\gdef\obiettivocorrente{}}
\bonusqformat{\textbf{\thequestion.}\enspace{\footnotesize\color{black!55}\textit{(facoltativo)}}\enspace\stampaobiettivo\hfill\colonnapunti{(+\totalbonuspoints\ bonus)}\gdef\obiettivocorrente{}}
\renewcommand{\questionlabel}{}
\renewcommand{\partlabel}{\thepartno)}
\renewcommand{\subpartlabel}{\thesubpart.}
\renewcommand{\choicelabel}{\thechoice)}
\pointpoints{punto}{punti}
\bonuspointpoints{punto bonus}{punti bonus}
\newcommand{\colonnapunti}[1]{{\small #1}}
\pointformat{\small(\thepoints)}
\bonuspointformat{\small(+\thepoints)}

% Sottopunti con interlinea maggiore (le soluzioni tornano all'interlinea normale)
\renewcommand{\partshook}{\linespread{\interlineaparts}\selectfont\setlength{\itemsep}{15pt}\setlength{\rightmargin}{2.2cm}}

% Numerazione degli esercizi: ogni fila riparte da 1 (esercizio 1, 2, ...);
% il contatore interno di exam prosegue da una fila all'altra, cosi' punti,
% totali e tabelle restano corretti. Usare l'ambiente "esercizi" al posto di "questions".
\newcounter{offsetfila}
\renewcommand{\thequestion}{\number\numexpr\value{question}-\value{offsetfila}\relax}
\newenvironment{esercizi}
{\begin{questions}\setcounter{question}{\value{offsetfila}}}
	{\end{questions}}

% ---------- Esercizi non spezzati tra due pagine ----------
%  \begin{esercizio} ... \end{esercizio} racchiude \obiettivo + \question + parts.
%  Se non ci sta, l'esercizio passa intero alla pagina successiva.
%  Nella versione con soluzioni l'esercizio puo' spezzarsi (i correttori sono lunghi).
\cancelspace   % le soluzioni nascoste non lasciano spazio ne' punti di interruzione
\newenvironment{esercizio}
{\par\ifprintanswers\else\addpenalty{-100}\begin{samepage}\fi}
	{\ifprintanswers\else\end{samepage}\fi\par\vspace{6pt}}

% ---------- Soluzioni e scelte multiple ----------
\renewcommand{\solutiontitle}{\noindent\textbf{Soluzione.}\enspace}
\unframedsolutions
\SolutionEmphasis{\color{blue!60!black}\linespread{\interlinea}\selectfont}
\CorrectChoiceEmphasis{\color{blue!60!black}\bfseries}
\checkboxchar{$\square$}
\checkedchar{$\blacksquare$}

% ---------- Tabella punteggi ----------
\hqword{\textbf{Esercizio}}
\hpword{\textbf{Punti massimi}}
\hsword{\textbf{Punti ottenuti}}
\htword{\textbf{Totale}}
\chqword{\textbf{Esercizio}}
\chpword{\textbf{Punti massimi}}
\chbpword{\textbf{Bonus}}
\chsword{\textbf{Punti ottenuti}}
\chtword{\textbf{Totale}}

% ---------- Macro utili ----------
\newcommand{\R}{\mathbb{R}}
\newcommand{\N}{\mathbb{N}}
\newcommand{\Z}{\mathbb{Z}}
\newcommand{\Q}{\mathbb{Q}}
\newcommand{\dd}{\,\mathrm{d}}

% ---------- Inizio e fine di ogni fila ----------
%  \iniziofila{A} ... \finefila   poi   \iniziofila{B} ... \finefila
%  \iniziofila: nuova pagina 1, intestazione scuola, titolo, nome, indicazioni.
%  \finefila:   tabella punteggi della sola fila, voto, firma.
\newcommand{\iniziofila}[1]{%
	\clearpage
	\renewcommand{\fila}{#1}%
	\setcounter{page}{1}%
	\setcounter{offsetfila}{\value{totalq@fila}}%
	\begin{center}
		{\Large \argomento}\\[2pt]
		{\small \indirizzo\ -- Classe \classe\ -- Fila \fila\ -- Durata: \durata}
	\end{center}
	\vspace{4pt}
	\noindent
	\begin{tabularx}{\textwidth}{@{}l X l p{3.2cm}@{}}
		\textbf{Nome e Cognome:} & \hrulefill & \textbf{Data:} & \hrulefill \\[8pt]
	\end{tabularx}
	\vspace{2pt}
	\noindent\fbox{\parbox{\dimexpr\textwidth-2\fboxsep-2\fboxrule}{\small\linespread{1.2}\selectfont
			\textbf{Indicazioni.} \indicazioni\
			Punteggio totale: \pointsinrange{fila#1} punti.}}
	\vspace{2pt}
	\begingradingrange{fila#1}}
\newcounter{totalq@fila}
\newcommand{\finefila}{%
	\setcounter{totalq@fila}{\value{question}}%
	\endgradingrange{fila\fila}%
	\vfill
	\begin{center}\small
		\iftabellabonus
		\partialcombinedgradetable{fila\fila}[h][questions]%
		\else
		\partialgradetable{fila\fila}[h][questions]%
		\fi
	\end{center}
	\vspace{6pt}
	\noindent{\small\textbf{Voto:} \rule{3cm}{0.4pt}\hfill
		\textbf{Firma del docente:} \rule{5cm}{0.4pt}}%
	\label{fine@fila@\fila}}

% =====================================================================
\begin{document}
	
	% =====================================================================
	% FILA A
	% =====================================================================
	\iniziofila{A}
	\begin{esercizi}
		
		% ---- ESERCIZIO 1 ----
		\begin{esercizio}
			\obiettivo{valutare trasformazioni equivalenti, soluzioni e propriet\`a delle equazioni}
			\question
			Stabilisci se ciascuna delle seguenti affermazioni \`e Vera (V) o Falsa (F), fornendo una breve motivazione concettuale.
			\begin{parts}
\part[1] Moltiplicare entrambi i membri di un'equazione intera per $x$ produce sempre un'equazione equivalente a quella iniziale.
\begin{solution}
\textbf{Falso.} Per esempio, $x-2=0$ ha $S=\{2\}$, mentre $x(x-2)=0$ ha $S=\{0;\,2\}$. Il fattore $x$ pu\`o introdurre la soluzione $x=0$.
\end{solution}
\part[1] Se, sostituendo $x=2$ in un'equazione, si ottiene $5=7$, allora $2$ non appartiene al suo insieme delle soluzioni.
\begin{solution}
\textbf{Vero.} Una soluzione deve rendere uguali i due membri. L'uguaglianza ottenuta \`e falsa, quindi $2\notin S$.
\end{solution}
\part[1] L'equazione $(x-2)^2=0$ ha una sola soluzione reale, pur essendo di secondo grado.
\begin{solution}
\textbf{Vero.} $(x-2)^2=(x-2)(x-2)$: entrambi i fattori si annullano soltanto per $x=2$. La radice ripetuta va scritta una sola volta nell'insieme $S=\{2\}$.
\end{solution}
\part[1] La presenza di coefficienti frazionari rende necessariamente frazionaria un'equazione.
\begin{solution}
\textbf{Falso.} Per esempio, $\frac{x}{2}+\frac13=1$ \`e intera: i denominatori sono numeri non nulli e non contengono l'incognita. Un coefficiente frazionario non rende l'equazione frazionaria.
\end{solution}
\end{parts}
		\end{esercizio}
		
		% ---- ESERCIZIO 2 ----
		\begin{esercizio}
			\obiettivo{semplificare e classificare equazioni intere riconducibili a forme lineari}
			\question
			Risolvi in $\R$ le seguenti equazioni intere, semplificando i denominatori numerici e sviluppando o riconoscendo i prodotti notevoli. Specifica se ciascuna \`e determinata, indeterminata o impossibile e scrivi l'insieme $S$ delle soluzioni.
			\begin{parts}
\part[2]
\[\frac{2x-1}{\frac12+\frac13}-\frac{x+2}{\frac34-\frac14}
=\frac{x-4}{1-\frac16}\]
\begin{solution}
I denominatori numerici sono $\frac56$, $\frac12$, $\frac56$, tutti non nulli; il dominio \`e $\R$.
\[
\frac65(2x-1)-2(x+2)=\frac65(x-4)
\]
Moltiplicando per $5$:
\[
6(2x-1)-10(x+2)=6(x-4)
\implies 2x-26=6x-24
\implies -4x=2.
\]
Determinata: $S=\left\{-\frac12\right\}$.
\end{solution}
\part[1]
\[(2x-1)^2-(x+2)(x-2)=3x(x-2)+2x+5\]
\begin{solution}
Sviluppando il quadrato e la differenza di quadrati:
\[
4x^2-4x+1-x^2+4=3x^2-6x+2x+5
\]
\[
3x^2-4x+5=3x^2-4x+5\implies 0=0.
\]
Indeterminata: $S=\R$.
\end{solution}
\part[1]
\[\frac{(3x-1)^2-(3x+1)(3x-1)}{2}=-3(x+1)+\frac52\]
\begin{solution}
Si pu\`o raccogliere $(3x-1)$ al numeratore:
\[
\frac{(3x-1)[(3x-1)-(3x+1)]}{2}=-3x-3+\frac52.
\]
Quindi $-3x+1=-3x-\frac12$, da cui $1=-\frac12$.
Impossibile: $S=\varnothing$.
\end{solution}
\part[3]
\[
\begin{aligned}
\left(2x-\frac13\right)^2+x^2\left(x^2+x-\frac5{12}\right)
&=\left(2x-\frac13\right)\left(2x+\frac13\right)\\
&\quad+\left(x^2+\frac{x}{2}-\frac13\right)^2.
\end{aligned}
\]
\begin{solution}
Sviluppando separatamente i due membri:
\[
\begin{aligned}
\mathrm{I}&=4x^2-\frac43x+\frac19+x^4+x^3-\frac5{12}x^2,\\
\mathrm{II}&=4x^2-\frac19+x^4+x^3-\frac5{12}x^2-\frac{x}{3}+\frac19.
\end{aligned}
\]
Si cancellano i termini di quarto, terzo e secondo grado:
\[
-\frac43x+\frac19=-\frac{x}{3}
\implies -x+\frac19=0.
\]
Determinata: $S=\left\{\frac19\right\}$.
\end{solution}
\end{parts}
		\end{esercizio}
		
		% ---- ESERCIZIO 3 ----
		\begin{esercizio}
			\obiettivo{ricavare variabili da formule letterali della Fisica}
			\question
			Ricava la variabile indicata tra parentesi da ciascuna delle seguenti formule della Fisica, considerando tutte le grandezze coinvolte come ben definite e non nulle.
			\begin{parts}
				\part[1] Legge del moto rettilineo uniforme: $s = s_0 + v \cdot t \qquad [v]$
				\begin{solution}
					$v = \dfrac{s - s_0}{t}$.
				\end{solution}
				\part[1] Prima legge di Ohm: $V = R \cdot I \qquad [R]$
				\begin{solution}
					$R = \dfrac{V}{I}$.
				\end{solution}
				\part[1] Pressione sul fondo di un recipiente: $P=P_0+dgh \qquad [h]$
				\begin{solution}
					$h = \dfrac{P-P_0}{dg}$.
				\end{solution}
			\end{parts}
		\end{esercizio}
		
		% ---- ESERCIZIO 4 ----
		\begin{esercizio}
			\obiettivo{risolvere equazioni di grado superiore mediante scomposizione e legge di annullamento}
			\question
			Risolvi in $\R$ le seguenti equazioni mediante scomposizione e legge di annullamento del prodotto. Scrivi l'insieme $S$ delle soluzioni.
			\begin{parts}
\part[1]
\[x^2-7x+8=2\]
\begin{solution}
Si porta il termine del secondo membro a sinistra:
\[
x^2-7x+6=0.
\]
Si cercano due numeri di somma $-7$ e prodotto $6$: sono $-1$ e $-6$.
\[
(x-1)(x-6)=0,\qquad S=\{1;\,6\}.
\]
\end{solution}
\part[2]
\[\frac{(2x-3)^2}{1-\frac14}-\frac{(x+1)^2}{\frac12-\frac16}=0\]
\begin{solution}
I denominatori sono $\frac34$ e $\frac13$. Moltiplicando per $3$:
\[
4(2x-3)^2-9(x+1)^2=0.
\]
Differenza di quadrati:
\[
[2(2x-3)-3(x+1)][2(2x-3)+3(x+1)]=0
\]
\[
(x-9)(7x-3)=0,\qquad S=\left\{\frac37;\,9\right\}.
\]
\end{solution}
\part[4]
\[x^5-2x^4-18x^3+36x^2+81x-162=0\]
\begin{solution}
Raggruppando le coppie di termini:
\[
x^4(x-2)-18x^2(x-2)+81(x-2)=0
\]
\[
(x-2)(x^4-18x^2+81)=(x-2)(x^2-9)^2=0.
\]
Il trinomio in $x^2$ \`e un quadrato perfetto. Scomponendo ancora:
\[
(x-2)(x-3)^2(x+3)^2=0,\qquad S=\{-3;\,2;\,3\}.
\]
\end{solution}
\end{parts}
		\end{esercizio}
		
		% ---- ESERCIZIO 5 ----
\begin{esercizio}
\obiettivo{modellizzare una spesa con sconto parziale e costo fisso}
\question
Una studentessa acquista quattro quaderni e tre raccoglitori. Ogni raccoglitore costa $3$ euro pi\`u di un quaderno. Il negozio applica ai soli raccoglitori uno sconto pari a un quinto del loro prezzo. Aggiungendo $2$ euro di spedizione, la studentessa paga complessivamente \num{41,20} euro.
\begin{parts}
\part[1] Scegli l'incognita e imposta un'equazione per determinare il prezzo di un quaderno.
\begin{solution}
Si pone $x$ uguale al prezzo, in euro, di un quaderno, con $x>0$.
Un raccoglitore costa $x+3$ euro; dopo lo sconto viene pagato $\frac45(x+3)$ euro.
\[
4x+3\cdot\frac45(x+3)+2=\num{41,20}.
\]
Lo sconto riguarda soltanto i raccoglitori: quaderni e spedizione si pagano interamente.
\end{solution}
\part[2] Calcola il prezzo di un quaderno e quello di un raccoglitore prima dello sconto.
\begin{solution}
Moltiplicando l'equazione per $5$:
\[
20x+12(x+3)+10=206
\]
\[
32x+46=206\implies32x=160\implies x=5.
\]
L'equazione ha $S=\{5\}$: un quaderno costa $5$ euro e un raccoglitore costa $5+3=8$ euro prima dello sconto.
Verifica della spesa:
\[
4\cdot5+3\cdot\frac45\cdot8+2=\num{41,20}.
\]
\end{solution}
\part[1] Quanto avrebbe speso in pi\`u senza lo sconto, mantenendo le stesse spese di spedizione?
\begin{solution}
Senza sconto avrebbe pagato $4\cdot5+3\cdot8+2=46$ euro.
La differenza \`e $46-\num{41,20}=\num{4,80}$ euro.
Equivalentemente, lo sconto totale \`e $3\cdot\frac15\cdot8=\num{4,80}$ euro.
\end{solution}
\end{parts}
\end{esercizio}
		
	\end{esercizi}
	\finefila
	
	% =====================================================================
	% FILA B
	% =====================================================================
	\iniziofila{B}
	\begin{esercizi}
		
		% ---- ESERCIZIO 1 ----
		\begin{esercizio}
			\obiettivo{valutare trasformazioni equivalenti, soluzioni e propriet\`a delle equazioni}
			\question
			Stabilisci se ciascuna delle seguenti affermazioni \`e Vera (V) o Falsa (F), fornendo una breve motivazione concettuale.
			\begin{parts}
\part[1] Dividendo entrambi i membri di $x^2=4x$ per $x$, si ottiene un'equazione equivalente a quella iniziale.
\begin{solution}
\textbf{Falso.} L'equazione iniziale diventa $x(x-4)=0$ e ha $S=\{0;\,4\}$. La divisione per $x$ richiede $x\ne0$ e fa perdere la soluzione $x=0$.
\end{solution}
\part[1] Le equazioni $2(x-1)=6$ e $x+3=7$ sono equivalenti, anche se hanno forme diverse.
\begin{solution}
\textbf{Vero.} Entrambe sono definite in $\R$ e hanno lo stesso insieme delle soluzioni, $S=\{4\}$. L'equivalenza dipende dalle soluzioni, non dall'aspetto delle equazioni.
\end{solution}
\part[1] Per la legge di annullamento del prodotto, una soluzione di $(x-1)(x+2)=0$ deve annullare contemporaneamente entrambi i fattori.
\begin{solution}
\textbf{Falso.} Basta che sia nullo almeno uno dei fattori: $x-1=0$ \emph{oppure} $x+2=0$. Dunque $S=\{-2;\,1\}$; nessun numero annulla qui entrambi i fattori contemporaneamente.
\end{solution}
\part[1] Il passaggio da $3x-5=2x+7$ a $3x-2x=7+5$ conserva l'insieme delle soluzioni.
\begin{solution}
\textbf{Vero.} Si sottrae $2x$ e si aggiunge $5$ a entrambi i membri: sono trasformazioni equivalenti. Entrambe le equazioni hanno $S=\{12\}$.
\end{solution}
\end{parts}
		\end{esercizio}
		
		% ---- ESERCIZIO 2 ----
		\begin{esercizio}
			\obiettivo{semplificare e classificare equazioni intere riconducibili a forme lineari}
			\question
			Risolvi in $\R$ le seguenti equazioni intere, semplificando i denominatori numerici e sviluppando o riconoscendo i prodotti notevoli. Specifica se ciascuna \`e determinata, indeterminata o impossibile e scrivi l'insieme $S$ delle soluzioni.
			\begin{parts}
\part[2]
\[\frac{2x+1}{\frac12+\frac14}-\frac{x-1}{\frac56-\frac13}
=\frac{x+2}{1-\frac14}\]
\begin{solution}
I denominatori numerici sono $\frac34$, $\frac12$, $\frac34$, tutti non nulli; il dominio \`e $\R$.
\[
\frac43(2x+1)-2(x-1)=\frac43(x+2).
\]
Moltiplicando per $3$:
\[
4(2x+1)-6(x-1)=4(x+2)
\implies 2x+10=4x+8
\implies -2x=-2.
\]
Determinata: $S=\{1\}$.
\end{solution}
\part[1]
\[\frac{(3x+1)^2-(3x-1)(3x+1)}{2}=3(x-1)+\frac52\]
\begin{solution}
Si pu\`o raccogliere $(3x+1)$ al numeratore:
\[
\frac{(3x+1)[(3x+1)-(3x-1)]}{2}=3x-3+\frac52.
\]
Quindi $3x+1=3x-\frac12$, da cui $1=-\frac12$.
Impossibile: $S=\varnothing$.
\end{solution}
\part[1]
\[(2x+1)^2-(x+3)(x-3)=3x(x+2)-2x+10\]
\begin{solution}
Sviluppando il quadrato e la differenza di quadrati:
\[
4x^2+4x+1-x^2+9=3x^2+6x-2x+10
\]
\[
3x^2+4x+10=3x^2+4x+10\implies 0=0.
\]
Indeterminata: $S=\R$.
\end{solution}
\part[3]
\[
\begin{aligned}
\left(3x+\frac13\right)^2+x^2\left(x^2-x-\frac5{12}\right)
&=\left(3x-\frac13\right)\left(3x+\frac13\right)\\
&\quad+\left(x^2-\frac{x}{2}-\frac13\right)^2.
\end{aligned}
\]
\begin{solution}
Sviluppando separatamente i due membri:
\[
\begin{aligned}
\mathrm{I}&=9x^2+2x+\frac19+x^4-x^3-\frac5{12}x^2,\\
\mathrm{II}&=9x^2-\frac19+x^4-x^3-\frac5{12}x^2+\frac{x}{3}+\frac19.
\end{aligned}
\]
Si cancellano i termini di quarto, terzo e secondo grado:
\[
2x+\frac19=\frac{x}{3}
\implies \frac53x+\frac19=0
\implies 15x+1=0.
\]
Determinata: $S=\left\{-\frac1{15}\right\}$.
\end{solution}
\end{parts}
		\end{esercizio}
		
		% ---- ESERCIZIO 3 ----
		\begin{esercizio}
			\obiettivo{ricavare variabili da formule letterali della Fisica}
			\question
			Ricava la variabile indicata tra parentesi da ciascuna delle seguenti formule della Fisica, considerando tutte le grandezze coinvolte come ben definite e non nulle.
			\begin{parts}
				\part[1] Pressione: $P = \dfrac{F}{S} \qquad [S]$
				\begin{solution}
					$S = \dfrac{F}{P}$.
				\end{solution}
				\part[1] Accelerazione media: $a = \dfrac{v - v_0}{t} \qquad [v]$
				\begin{solution}
					$v = v_0 + a t$.
				\end{solution}
				\part[1] Pressione sul fondo di un recipiente: $P=P_0+dgh \qquad [h]$
				\begin{solution}
					$h = \dfrac{P-P_0}{dg}$.
				\end{solution}
			\end{parts}
		\end{esercizio}
		
		% ---- ESERCIZIO 4 ----
		\begin{esercizio}
			\obiettivo{risolvere equazioni di grado superiore mediante scomposizione e legge di annullamento}
			\question
			Risolvi in $\R$ le seguenti equazioni mediante scomposizione e legge di annullamento del prodotto. Scrivi l'insieme $S$ delle soluzioni.
			\begin{parts}
\part[1]
\[x^2+5x+3=-3\]
\begin{solution}
Si porta il termine del secondo membro a sinistra:
\[
x^2+5x+6=0.
\]
Si cercano due numeri di somma $5$ e prodotto $6$: sono $2$ e $3$.
\[
(x+2)(x+3)=0,\qquad S=\{-3;\,-2\}.
\]
\end{solution}
\part[2]
\[\frac{(2x+1)^2}{1-\frac14}-\frac{(x-2)^2}{\frac12-\frac16}=0\]
\begin{solution}
I denominatori sono $\frac34$ e $\frac13$. Moltiplicando per $3$:
\[
4(2x+1)^2-9(x-2)^2=0.
\]
Differenza di quadrati:
\[
[2(2x+1)-3(x-2)][2(2x+1)+3(x-2)]=0
\]
\[
(x+8)(7x-4)=0,\qquad S=\left\{-8;\,\frac47\right\}.
\]
\end{solution}
\part[4]
\[x^5+3x^4-8x^3-24x^2+16x+48=0\]
\begin{solution}
Raggruppando le coppie di termini:
\[
x^4(x+3)-8x^2(x+3)+16(x+3)=0
\]
\[
(x+3)(x^4-8x^2+16)=(x+3)(x^2-4)^2=0.
\]
Il trinomio in $x^2$ \`e un quadrato perfetto. Scomponendo ancora:
\[
(x+3)(x-2)^2(x+2)^2=0,\qquad S=\{-3;\,-2;\,2\}.
\]
\end{solution}
\end{parts}
		\end{esercizio}
		
		% ---- ESERCIZIO 5 ----
\begin{esercizio}
\obiettivo{modellizzare una spesa con sconto parziale e costo fisso}
\question
Uno studente acquista tre quaderni e quattro raccoglitori. Ogni raccoglitore costa $2$ euro pi\`u di un quaderno. Il negozio applica ai soli raccoglitori uno sconto pari a un quarto del loro prezzo. Aggiungendo $3$ euro di spedizione, lo studente paga complessivamente $45$ euro.
\begin{parts}
\part[1] Scegli l'incognita e imposta un'equazione per determinare il prezzo di un quaderno.
\begin{solution}
Si pone $x$ uguale al prezzo, in euro, di un quaderno, con $x>0$.
Un raccoglitore costa $x+2$ euro; dopo lo sconto viene pagato $\frac34(x+2)$ euro.
\[
3x+4\cdot\frac34(x+2)+3=45.
\]
Lo sconto riguarda soltanto i raccoglitori: quaderni e spedizione si pagano interamente.
\end{solution}
\part[2] Calcola il prezzo di un quaderno e quello di un raccoglitore prima dello sconto.
\begin{solution}
Semplificando:
\[
3x+3(x+2)+3=45
\]
\[
6x+9=45\implies6x=36\implies x=6.
\]
L'equazione ha $S=\{6\}$: un quaderno costa $6$ euro e un raccoglitore costa $6+2=8$ euro prima dello sconto.
Verifica della spesa:
\[
3\cdot6+4\cdot\frac34\cdot8+3=45.
\]
\end{solution}
\part[1] Quanto avrebbe speso in pi\`u senza lo sconto, mantenendo le stesse spese di spedizione?
\begin{solution}
Senza sconto avrebbe pagato $3\cdot6+4\cdot8+3=53$ euro.
La differenza \`e $53-45=8$ euro.
Equivalentemente, lo sconto totale \`e $4\cdot\frac14\cdot8=8$ euro.
\end{solution}
\end{parts}
\end{esercizio}
		
	\end{esercizi}
	\finefila
	
\end{document}
